\documentclass[11pt,letterpaper]{article}
\usepackage[T1]{fontenc}
\usepackage{lmodern,microtype}
\usepackage[margin=0.78in,headheight=15pt,headsep=15pt]{geometry}
\usepackage{mathtools,amssymb}
\usepackage{xcolor,booktabs,array,tikz}
\usepackage[most]{tcolorbox}
\usepackage{enumitem,fancyhdr,needspace,hyperref}
\definecolor{ink}{HTML}{162332}
\definecolor{navy}{HTML}{17324D}
\definecolor{teal}{HTML}{167D82}
\definecolor{mist}{HTML}{EDF5F4}
\definecolor{muted}{HTML}{5D6B78}
\color{ink}
\hypersetup{colorlinks=true,linkcolor=teal,urlcolor=teal,pdfauthor={Igor Rivin}}
\setlength{\parindent}{0pt}
\setlength{\parskip}{4pt}
\setlength{\abovedisplayskip}{7pt}
\setlength{\belowdisplayskip}{7pt}
\setlist[enumerate]{label=(\alph*),leftmargin=1.8em,itemsep=3pt,topsep=4pt}
\newcommand{\Z}{\mathbb Z}
\newcommand{\F}{\mathbb F}
\newcommand{\one}{\mathbf{1}}
\newcommand{\id}{\operatorname{id}}
\newcommand{\pts}[1]{\hfill\textcolor{muted}{\small[#1 points]}}
\newcommand{\partpts}[1]{\unskip\nobreak\hspace{.4em}\textcolor{muted}{\small(#1)}}
\newcommand{\edition}{\ifsolutions Instructor solutions\else Student handout\fi}
\newcounter{problem}
\newcommand{\problem}[2]{\par\addvspace{11pt}\Needspace{6\baselineskip}%
  \stepcounter{problem}{\bfseries\color{navy}\theproblem.\ #1}\pts{#2}\par\nobreak}
\newcommand{\topic}[1]{\par\addvspace{9pt}\Needspace{10\baselineskip}%
  {\large\bfseries\color{teal}#1}\par\nobreak}
\newcommand{\studentbreak}{\ifsolutions\else\newpage\fi}
\newcommand{\challenge}[1]{\par\addvspace{12pt}\Needspace{7\baselineskip}%
  {\bfseries\color{teal}Optional challenge: #1}\pts{10 extra}\par\nobreak}
\pagestyle{fancy}
\fancyhf{}
\fancyhead[L]{\small\color{muted}\setname}
\fancyhead[R]{\small\color{muted}Patterns, Proofs, and Primes}
\fancyfoot[L]{\footnotesize\color{muted}\edition}
\fancyfoot[R]{\small\color{teal}\thepage}
\renewcommand{\headrulewidth}{0.3pt}
\NewDocumentEnvironment{solution}{+b}{\ifsolutions
  \begin{tcolorbox}[enhanced,breakable,colback=mist!45,colframe=teal!45,
    boxrule=.5pt,arc=1mm,left=3mm,right=3mm,top=2mm,bottom=2mm,
    before skip=7pt,after skip=7pt,title={Solution},
    title after break={Solution (continued)},colbacktitle=mist,
    coltitle=teal,fonttitle=\small\bfseries]
  \small #1
  \end{tcolorbox}\fi}{}
\newcommand{\handouttitle}[3]{%
  {\LARGE\bfseries\color{navy}\setname}\hfill{\small\color{teal}\edition}\\[4pt]
  {\large #1}\\[4pt]{\small\color{muted}#2}\par
  \medskip\textbf{Preparation.} #3\par
  \ifsolutions\else\medskip
  \textbf{Name:}\ \rule{.46\linewidth}{.4pt}\hfill
  \textbf{Due:}\ \rule{.20\linewidth}{.4pt}\par\fi}

\newcommand{\Q}{\mathbb Q}
\newcommand{\setname}{Valuations and $p$-adic distance}
\hypersetup{pdftitle={Number Theory: Valuations, Primes, and p-adic Distance},
  pdfsubject={Factorial and binomial valuations, Chebyshev estimates, and introductory p-adic analysis}}
\begin{document}
\handouttitle{From factorials and prime counting to a new notion of closeness}
{7 problems \textbar\ 100 points + 10 challenge points
\textbar\ Suggested time: 3--4 hours}
{The lecture on Chebyshev's estimate and factorial/binomial valuations;
course packet Chapters 7--8. No prior study of metric spaces is assumed.}

\textbf{Instructions.} Write solutions on separate paper and justify your
answers. You may use unique factorization, the binomial theorem, and
integrality of binomial coefficients. You may use results of earlier problems
even if you have not proved them. All logarithms are natural.

\textbf{Notation.} For a prime $p$ and a nonzero integer $m$, $v_p(m)$ is
the exponent of $p$ in $|m|$. Set $v_p(0)=+\infty$, with
$+\infty+t=+\infty$ for finite $t$. Recall that
$\pi(x)=\#\{p\le x:p\text{ prime}\}$.

\ifsolutions
\textbf{Grading guidance.} Subpart scores total 100; the challenge contributes
10 extra points. Problems 1--3 review and apply the lecture, and Problems
4--7 introduce $p$-adic distance and convergence entirely on $\Q$.
Accept equivalent proofs and carry forward arithmetic errors when the later
method is correct. Do not require a construction of $\Q_p$ or any general
theory of metric spaces.
\fi

\topic{I. Factorials, binomial coefficients, and primes}

\problem{Count the powers of a prime}{12}
\begin{enumerate}
 \item Prove Legendre's formula by counting multiples of $p,p^2,\ldots$:
 \[
 v_p(n!)=\sum_{j\ge1}\left\lfloor\frac{n}{p^j}\right\rfloor.
 \]
 Explain why multiples of $p^2$ must be counted more than once and why
 the sum has only finitely many nonzero terms. \partpts{4}
 \item Compute $v_2(50!)$, $v_3(50!)$, and $v_5(50!)$.
 How many trailing zeros does $50!$ have in base ten? \partpts{4}
 \item Show that $v_p((pn)!)=n+v_p(n!)$. Deduce that, for $r\ge1$,
 \[
 v_p((p^r)!)=1+p+\cdots+p^{r-1}=\frac{p^r-1}{p-1}.
 \]
 \partpts{4}
\end{enumerate}
\begin{solution}
Each integer $m\le n$ contributes one copy of $p$ for each $j\ge1$ with
$p^j\mid m$. Thus summing these contributions first over $m$, then over
$j$, gives the formula. A multiple of $p^2$ contributes at least two copies;
no $p^j>n$ contributes.
The requested values are
\[
 v_2(50!)=25+12+6+3+1=47,\quad
 v_3(50!)=16+5+1=22,\quad v_5(50!)=10+2=12.
\]
Since $10^k\mid50!$ exactly when both $2^k$ and $5^k$ divide it, the
number of trailing zeros is $\min(47,12)=12$.
Separating the first term of Legendre's formula gives
$v_p((pn)!)=n+\sum_{j\ge1}\lfloor n/p^j\rfloor$.
Apply this repeatedly, starting with $v_p(1!)=0$, to get the geometric sum.
\end{solution}

\problem{Why the middle of Pascal's triangle fits inside an lcm}{16}
Let $n\ge1$ and $B_n=\binom{2n}{n}$.
\begin{enumerate}
 \item Use Problem 1 to show that
 \[
 v_p(B_n)=\sum_{j\ge1}\left(
 \left\lfloor\frac{2n}{p^j}\right\rfloor
 -2\left\lfloor\frac{n}{p^j}\right\rfloor\right).
 \]
 Prove that every summand is either $0$ or $1$.
 \emph{Hint:} Write $n/p^j=q+t$, with $q\in\Z$ and $0\le t<1$.
 \partpts{5}
 \item Compute $v_2\binom{20}{10}$, $v_3\binom{20}{10}$, and
 $v_5\binom{20}{10}$ without expanding the binomial coefficient.
 \partpts{3}
 \item Put $L_n=\operatorname{lcm}(1,2,\ldots,2n)$. Determine $v_p(L_n)$
 and prove that $B_n\mid L_n$ by comparing valuations for every prime.
 \partpts{5}
 \item Prove that if $n<p\le2n$, then $v_p(B_n)=1$.
 \partpts{3}
\end{enumerate}
\begin{solution}
Subtract twice the valuation of $n!$ from the valuation of $(2n)!$.
Writing $n/p^j=q+t$ gives
$\lfloor2(q+t)\rfloor-2\lfloor q+t\rfloor=\lfloor2t\rfloor\in\{0,1\}$.
For $n=10$, Legendre's formula yields
\[
 v_2(B_{10})=18-2\cdot8=2,\quad
 v_3(B_{10})=8-2\cdot4=0,\quad
 v_5(B_{10})=4-2\cdot2=0.
\]
The largest power of $p$ occurring among $1,\ldots,2n$ is $p^e$, where
$e=\lfloor\log(2n)/\log p\rfloor$ (this is $0$ when $p>2n$).
Thus $v_p(L_n)=e$. There are at most $e$ nonzero summands in part (a),
each at most $1$, so $v_p(B_n)\le e$. Unique factorization proves divisibility.
Finally, $n<p\le2n$ gives $\lfloor2n/p\rfloor=1$ and
$\lfloor n/p\rfloor=0$. Also $p^2>2n$: since $p\ge n+1$,
$p^2\ge(n+1)^2>2n$. All later terms vanish.
\end{solution}

\studentbreak
\problem{Recover the scale of prime counting}{16}
\begin{enumerate}
 \item Explain why the central coefficient is largest in row $2n$, and prove
 $\frac{4^n}{2n+1}\le B_n<4^n$.
 \emph{Hint:} Compare consecutive coefficients and sum the row. \partpts{4}
 \item Using Problem 2, prove
 \[
 \log B_n\le\pi(2n)\log(2n),\qquad
 \pi(2n)\ge\frac{2n\log2-\log(2n+1)}{\log(2n)}.
 \]
 Explain why monotonicity of $\pi$ gives a constant $c>0$ such that
 $\pi(x)\ge c\,x/\log x$ for all sufficiently large real $x$.
 \partpts{6}
 \item You may take the following lecture estimate as given:
 $\vartheta(x)=\sum_{p\le x}\log p\le Ax$ for $x\ge2$, where $A>0$
 is a fixed constant. Split the primes at $\sqrt{x}$ to prove
 $\pi(x)\le\sqrt{x}+\frac{2Ax}{\log x}$ for $x\ge4$.
 Deduce $\pi(x)\le Cx/\log x$ for some fixed $C$ and sufficiently large $x$.
 \partpts{4}
 \item Do these bounds prove $\pi(x)\sim x/\log x$? Explain the distinction
 between $\asymp$ and $\sim$. \partpts{2}
\end{enumerate}
\begin{solution}
The ratio of consecutive coefficients is
$\binom{2n}{k+1}/\binom{2n}{k}=(2n-k)/(k+1)$, so they increase up to
$k=n$ and decrease afterward. Their sum is $4^n$, giving the lower
bound; positive terms other than $B_n$ make the upper bound strict.
For each $p\le2n$, Problem 2 gives
$v_p(B_n)\log p\le\lfloor\log(2n)/\log p\rfloor\log p\le\log(2n)$.
Sum over these primes and combine with part (a).
For real $x$, take $n=\lfloor x/2\rfloor$. Then $2n\le x$,
$2n/x\to1$, and $\log(2n)/\log x\to1$.
The resulting lower bound, divided by $x/\log x$, tends to $\log2$;
thus $c=(\log2)/2$ works for sufficiently large $x$.
Every prime in $(\sqrt{x},x]$ contributes at least $\tfrac12\log x$ to
$\vartheta(x)$, so
$(\pi(x)-\pi(\sqrt{x}))\tfrac12\log x\le Ax$.
Use $\pi(\sqrt{x})\le\sqrt{x}$ and then
$\log x/\sqrt{x}\to0$; for example $C=2A+1$ eventually works.
The two bounds prove $\pi(x)\asymp x/\log x$, meaning that the ratio is
bounded above and below by positive constants for large $x$.
They do not prove $\sim$, which requires the ratio to tend to $1$.
\end{solution}

\topic{II. Divisibility becomes a notion of distance}

\problem{Extend the valuation to rational numbers}{14}
For $a,b\in\Z\setminus\{0\}$, define
$v_p(a/b)=v_p(a)-v_p(b)$; keep $v_p(0)=+\infty$.
\begin{enumerate}
 \item Prove that this definition is independent of the representation
 of a rational number as a fraction. \partpts{4}
 \item For nonzero $x,y\in\Q$, prove $v_p(xy)=v_p(x)+v_p(y)$ and
 $v_p(x^{-1})=-v_p(x)$. Compute
 $v_2(12/25)$, $v_5(12/25)$, and $v_3(9/14)$. \partpts{4}
 \item Prove, for all $x,y\in\Q$, that
 \[
 v_p(x+y)\ge\min\{v_p(x),v_p(y)\}.
 \]
 Show that equality holds when $v_p(x)\ne v_p(y)$.
 Give nonzero $x,y$ for which the inequality is strict and $x+y\ne0$.
 \emph{Hint:} First work with integers, factor out the smaller power of $p$,
 then use a common denominator. \partpts{6}
\end{enumerate}
\begin{solution}
If $a/b=c/d$, then $ad=bc$, so
$v_p(a)+v_p(d)=v_p(b)+v_p(c)$ and the two definitions agree.
Multiplication and inversion follow by adding or subtracting the integer
valuations of numerators and denominators. The numerical answers are
$2,-2,2$, respectively.
For nonzero integers $a,b$, let $r=\min(v_p(a),v_p(b))$ and write
$a=p^ru$, $b=p^rv$. Then $p^r\mid a+b$.
If the valuations differ, exactly one of $u,v$ is divisible by $p$, so
$u+v$ is not divisible by $p$, proving equality.
For rational $x=a/d$, $y=b/d$, with a common nonzero integer denominator,
subtract $v_p(d)$ from the integer inequality; the equality statement
also follows. Zero arguments are immediate under our convention.
For strict inequality, take $x=1$, $y=p-1$: both have valuation $0$,
but $x+y=p$ has valuation $1$.
\end{solution}

\problem{The ultrametric inequality}{14}
Define the \emph{$p$-adic absolute value} by $|0|_p=0$ and
$|x|_p=p^{-v_p(x)}$ for $x\ne0$. Define the \emph{$p$-adic distance}
by $d_p(x,y)=|x-y|_p$.
\begin{enumerate}
 \item Prove $|xy|_p=|x|_p|y|_p$ and the \emph{ultrametric inequality}
 \[
 |x+y|_p\le\max\{|x|_p,|y|_p\}.
 \]
 When $|x|_p\ne|y|_p$, prove that equality holds. \partpts{6}
 \item Verify the defining properties of a metric: $d_p(x,y)\ge0$,
 equality to $0$ exactly when $x=y$, symmetry, and the triangle inequality.
 In fact, prove the stronger statement
 $d_p(x,z)\le\max\{d_p(x,y),d_p(y,z)\}$.
 \partpts{5}
 \item Compute $|81|_3$, $|1/27|_3$, and $|10|_3$.
 Which is closer to $0$ in $d_3$: $81$ or $1/27$?
 Compare with ordinary absolute value. \partpts{3}
\end{enumerate}
\begin{solution}
For nonzero factors, exponentiate the product identity from Problem 4;
zero factors also satisfy it. Since $t\mapsto p^{-t}$ is decreasing,
the valuation inequality becomes the ultrametric inequality, and unequal
valuations give equality. The case $x+y=0$ follows from $|0|_p=0$.
Nonnegativity and definiteness follow from the definition. Symmetry follows
from $v_p(-u)=v_p(u)$. Apply the ultrametric inequality to
$x-z=(x-y)+(y-z)$. Because $\max(a,b)\le a+b$ for nonnegative $a,b$,
the ordinary triangle inequality follows too.
The three absolute values are $3^{-4}=1/81$, $3^3=27$, and $1$.
Thus $81$ is closer to $0$ in $d_3$, whereas $1/27$ is closer in ordinary distance.
\end{solution}

\studentbreak
\problem{Closeness means agreement modulo a prime power}{12}
\begin{enumerate}
 \item For integers $a,b$ and $k\ge0$, prove
 \[
 d_p(a,b)\le p^{-k}\quad\Longleftrightarrow\quad a\equiv b\pmod{p^k}.
 \]
 \partpts{4}
 \item Compute $d_5(7,32)$, $d_5(7,132)$, and $d_5(7,8)$.
 Which of $32,132,8$ is closest to $7$ in this metric? \partpts{3}
 \item List all integers $b$ with $0\le b<125$ and $d_5(7,b)\le5^{-2}$.
 Explain why the full set of integers satisfying this distance bound
 forms a congruence class. \partpts{3}
 \item Show that $|m|_p\le1$ for every integer $m$.
 Must this hold for every rational number? Give a reason or counterexample.
 \partpts{2}
\end{enumerate}
\begin{solution}
For $a\ne b$, the distance bound is equivalent to $v_p(a-b)\ge k$,
which is exactly $p^k\mid a-b$. When $a=b$, both assertions hold.
The distances are $1/25$, $1/125$, and $1$, respectively, so $132$ is
closest. The integers in part (c) are $7,32,57,82,107$: all are $7$ modulo
$25$. The criterion asks for divisibility of a difference, which allows
arbitrarily large ordinary differences.
Nonzero integers have nonnegative valuation, so their absolute values
are at most $1$; zero also satisfies the bound. For a rational counterexample,
$|1/p|_p=p>1$.
\end{solution}

\topic{III. First examples of $p$-adic convergence}

\textbf{Definition.} A sequence $(x_n)$ of rational numbers
\emph{converges $p$-adically} to $L\in\Q$ if, for every real $\varepsilon>0$,
there is an integer $N$ such that $|x_n-L|_p<\varepsilon$ whenever $n\ge N$.
Equivalently, the ordinary real sequence $|x_n-L|_p$ tends to $0$.
For integer $x_n,L$, Problem 6 says this happens exactly when, for every
$k\ge1$, eventually $x_n\equiv L\pmod{p^k}$.

\problem{Large integers can converge to zero}{16}
\begin{enumerate}
 \item Prove $p^n\to0$ $p$-adically. Use Problem 1 to prove $n!\to0$
 $p$-adically for every fixed prime $p$. In each case give a bound on
 $|x_n|_p$ that tends to $0$. \partpts{5}
 \item Let $S_n=1+p+p^2+\cdots+p^n$. Use the finite geometric-sum identity
 to prove
 \[
 S_n\longrightarrow\frac{1}{1-p}\quad\text{$p$-adically}.
 \]
 Compute the exact distance to the proposed limit. Specialize to $p=2$:
 what is the limit of $1+2+4+\cdots+2^n$? \partpts{5}
 \item Let $T_n=\sum_{j=1}^n j!$. For any $m>n$, use the ultrametric
 inequality repeatedly to show
 \[
 |T_m-T_n|_p\le p^{-\lfloor(n+1)/p\rfloor}.
 \]
 Deduce that $(T_n)$ is \emph{Cauchy}: for every $\varepsilon>0$, there
 is $N$ such that $|T_m-T_n|_p<\varepsilon$ for all $m,n\ge N$.
 You are not asked to find a rational limit. \partpts{4}
 \item Prove that a sequence of rational numbers has at most one rational
 $p$-adic limit. \emph{Hint:} Apply the ultrametric inequality to
 $L-M=(L-x_n)+(x_n-M)$. \partpts{2}
\end{enumerate}
\begin{solution}
We have $|p^n|_p=p^{-n}\to0$. Legendre's formula gives
$v_p(n!)\ge\lfloor n/p\rfloor$, so $|n!|_p\le p^{-\lfloor n/p\rfloor}\to0$.
The geometric-sum identity gives
\[
 S_n-\frac{1}{1-p}=-\frac{p^{n+1}}{1-p}.
\]
Since $p\nmid(1-p)$, the distance is exactly $p^{-(n+1)}$, which tends
to $0$. For $p=2$ the limit is $-1$, with
$S_n=2^{n+1}-1\equiv-1\pmod{2^k}$ as soon as $n+1\ge k$.
For $m>n$, the finite-sum ultrametric inequality gives
\[
 |T_m-T_n|_p\le\max_{n+1\le j\le m}|j!|_p
 \le p^{-\lfloor(n+1)/p\rfloor}.
\]
Choose $N$ making the last expression less than $\varepsilon$ for
$n\ge N$. Symmetry covers $m<n$, and $m=n$ gives distance $0$.
This proves the Cauchy assertion; it does not by itself establish a limit
in $\Q$. The completion of $\Q$ for this metric is denoted $\Q_p$;
its construction lies beyond this assignment.
If $L$ and $M$ were both limits, then
$|L-M|_p\le\max(|L-x_n|_p,|x_n-M|_p)\to0$.
Definiteness forces $L=M$.
\end{solution}

\challenge{the geometry of an ultrametric}
\begin{enumerate}
 \item For three distinct rational numbers, prove that the largest of
 their three pairwise $p$-adic distances occurs at least twice. Thus every
 triangle is isosceles, with its two largest sides equal (equilateral
 triangles are allowed). \partpts{4}
 \item For a real $r>0$, define the closed ball
 $B(a,r)=\{x\in\Q:d_p(x,a)\le r\}$. Prove that if $b\in B(a,r)$,
 then $B(b,r)=B(a,r)$. \partpts{3}
 \item Prove that two closed balls are either disjoint or one is contained
 in the other. \emph{Hint:} If the balls intersect, choose a point in their
 intersection and apply part (b). \partpts{3}
\end{enumerate}
\begin{solution}
If one distance were strictly larger than the other two, the ultrametric
inequality applied to that pair through the third point would contradict it.
If $b\in B(a,r)$ and $x\in B(a,r)$, then
$d_p(x,b)\le\max(d_p(x,a),d_p(a,b))\le r$.
This proves $B(a,r)\subseteq B(b,r)$; interchange $a,b$ for the reverse inclusion.
If $B(a,r)$ and $B(b,s)$ contain a common point $c$, part (b) rewrites
them as $B(c,r)$ and $B(c,s)$. Assuming $r\le s$, the first is contained
in the second.
\end{solution}
\end{document}
